\chapter{执行摘要与阅读指南}

\section{项目定义与问题域}

本章用于给出整份报告的总导读。本文将 Pebble 定义为一个以关系上下文为核心的 AI 情绪防御与沟通辅助系统。它面向处在职场、家庭或亲密关系压力中的用户，目标是在高压互动发生时提供可持续的理解、练习与自我调节支持，而非只输出一次性的情绪标签或单轮回复建议。

从问题域看，Pebble 关注的是同一句话在不同关系结构、互动历史与当前情境中会呈现出不同的风险含义，因此系统需要同时处理语言表层、关系背景、应对策略与用户状态恢复四类信息。基于这一判断，Pebble 的产品主线不应被理解为单点问答工具，而应被理解为围绕复杂人际互动展开的连续支持链路。

\section{核心能力闭环}

Pebble 当前适合采用“四段闭环”的方式概括其能力结构。

\begin{itemize}
  \item 对话解码：用户输入一句令自己不适、困惑或难以回应的话语，系统输出表层语义、潜台词与候选回复方案。
  \item 关系陪练：用户在场景化对话中持续练习回应方式，并从系统获得即时反馈与策略提示。
  \item 呼吸稳态：用户在情绪负荷过高时进入短时稳定化流程，以降低后续判断和表达的失真风险。
  \item 关系档案与长期洞察：系统持续记录关系对象、互动背景和练习结果，使分析与建议能够逐步积累上下文。
\end{itemize}

为了让读者在进入后文章节前先建立稳定的全局心智模型，图~\ref{fig:pebble-overview-loop} 将 Pebble 的导读层压缩为一张总览图。该图强调 Pebble 的价值来自四段能力的闭环，而非某一个单点模块的孤立表现。

\begin{figure}[H]
  \centering
  \begin{tikzpicture}[
  node distance=1.7cm and 1.2cm,
  every node/.style={font=\small}
]
  \node[reportprocess] (decode) {对话解码\\识别潜台词与风险};
  \node[reportprocess, right=of decode] (practice) {关系陪练\\形成回应策略};
  \node[reportprocess, below=of practice] (archive) {关系档案\\沉淀长期上下文};
  \node[reportprocess, left=of archive] (breathing) {呼吸稳态\\降低情绪失真};
  \node[reportdecision, above right=0.85cm and 1.8cm of decode] (context) {关系上下文\\驱动全局解释};

  \draw[reportarrow] (decode) -- (practice);
  \draw[reportarrow] (practice) -- (archive);
  \draw[reportarrow] (archive) -- (breathing);
  \draw[reportarrow] (breathing) -- (decode);
  \draw[reportarrow] (context) -- ($(decode)!0.5!(practice)$);
  \draw[reportarrow] (context) -- ($(practice)!0.5!(archive)$);
  \draw[reportarrow] (context) -- ($(archive)!0.5!(breathing)$);
  \draw[reportarrow] (context) -- ($(breathing)!0.5!(decode)$);
\end{tikzpicture}

  \caption{Pebble 总览闭环，突出关系上下文如何统摄解码、陪练、稳态恢复与长期记忆。}
  \label{fig:pebble-overview-loop}
\end{figure}

这一闭环意味着 Pebble 的核心价值来自能力之间的联动。解码模块负责帮助用户理解当下对话；陪练模块负责把理解转化为可重复训练的表达能力；呼吸模块负责在高压时刻提供最小可用的生理稳定支持；关系档案负责把一次次局部互动沉淀为长期记忆，使系统能够在后续分析中引入更细粒度的情境差异。

\section{证据边界与当前定位}

本报告后续章节会同时使用两类信息来源。

\begin{itemize}
  \item design intent：用于说明 Pebble 试图建立什么样的产品结构、能力边界与用户体验。
  \item 已观察事实：用于说明哪些系统能力已经得到直接观察。
\end{itemize}

基于这一区分，本文后续论述会把系统愿景、已验证能力与交付边界明确拆开，以保证判断强度与证据强度保持一致。

与此同时，第一章不承担展开实现细节的职责。接口字段、认证路径、关系数据模型、测试矩阵、迁移历史与当前阻塞等内容会在后续章节分别展开，其中与真实实现状态直接相关的信息默认集中进入“项目进度”章节。本章只保留理解整份报告所必需的最小事实，以避免导读层被工程细节淹没。

\section{阅读建议}

如果读者首先关心 Pebble 要解决什么问题，以及它为何需要围绕关系上下文组织能力，建议优先阅读第二章与第三章。第二章会明确项目主张、目标边界与读者应如何理解 Pebble 的系统使命；第三章会展开背景、冲突与上下文，说明这一类产品为何不能仅依赖单句文本分析。

如果读者首先关心系统如何工作，建议继续阅读第四章至第八章。第四章负责定义关键术语和前置概念，第五章概括核心设计原则，第六章和第七章分别展开流程与系统资源设计，第八章说明示例、验证与可接受结果。若读者关心当前真实进展、里程碑与风险，则应直接阅读第九章和第十章。

读完整份报告后，读者应形成以下三个稳定判断。

\begin{itemize}
  \item Pebble 的系统边界建立在“关系上下文驱动的沟通支持”之上。
  \item Pebble 的产品能力应以解码、练习、稳态恢复与长期记忆的联动来理解。
  \item Pebble 已经具备进入架构与产品设计讨论的系统骨架，但实现成熟度、验证深度与交付边界仍需要在后文按证据层逐项区分。
\end{itemize}
